\section{Lean formalization plan}
\label{sec:lean-plan}

\subsection{Formalization status}

There are two distinct layers.

\begin{enumerate}
\item The \emph{semantic algebra layer} is already kernel-checked in the
  inspected ProveIt tree.  It contains the field-generic factorization
  theorems for degrees at most four, the typed radical-expression language,
  the semantic radical closure, explicit rational Abel--Ruffini examples, and
  the generic fixed-field construction.
\item The \emph{iPRA arithmetization layer} specified in
  Sections~\ref{sec:ipra-framework} and~\ref{sec:internalization} is not yet a
  completed Lean development.  The present archive gives its definitions,
  theorem interfaces, proof decomposition, and an exact finite checker, but
  it does not claim an end-to-end kernel-checked theorem of type
  \code{PRAAx derives AbelRuffini}.
\end{enumerate}

This distinction matters.  Importing a classical field-theory theorem into
Lean proves the mathematics, but it does not by itself prove that a derivation
exists in the weak object theory iPRA.  The missing bridge must internalize the
finite codes and the correctness proofs of their primitive-recursive
checkers.

\subsection{Existing files to reuse}

The following imports should remain the semantic reference implementation.

\begin{longtable}{p{0.34\textwidth}p{0.58\textwidth}}
\toprule
Existing module & Reused theorem families\\
\midrule
\endhead
\code{PolynomialFormulas.Basic} &
\code{linear_eq_zero_iff}, \code{quadratic_formula_factorization},
\code{quadratic_eq_zero_iff}.\\
\code{PolynomialFormulas.Cubic} &
\code{depress_monic_cubic}, \code{cardano_depressed},
\code{cardano_factorization}, \code{cubic_factorization},
\code{cubic_eq_zero_iff}.\\
\code{PolynomialFormulas.Quartic} &
\code{depress_monic_quartic}, \code{ferrari_factorization},
\code{solveDepressedQuartic_factorization},
\code{quartic_factorization}, \code{quartic_eq_zero_iff}.\\
\code{PolynomialFormulas.AbelRuffini} &
\code{RadicalExpression},
\code{nonempty_iff_mem_solvableByRad}, and the explicit rational quintic
obstruction.\\
\code{PolynomialFormulas.GenericAbelRuffini} &
The rational-function-field model, root orbit, splitting-field construction,
Galois equivalence with a symmetric group, and rootwise obstruction.\\
\code{ConciseIPRA.NativeBiInterpretation} &
\code{native_iff_pra}, \code{pra_iff_native}, and the two round-trip
translation theorems.\\
\bottomrule
\end{longtable}

The file \code{PolynomialFormulas.Audit} already asks Lean to print the axioms
of the main semantic theorems.  The new development should add a separate
audit for the iPRA endpoints and fail continuous integration if any endpoint
contains \code{sorryAx} or an unintended axiom.

\subsection{Proposed module tree}

A practical directory is
\[
 \code{Logic/PrimitiveRecursiveArithmetic/Applications/AbelRuffini/}.
\]
The following modules give a dependency order with no cycles.

\begin{description}[style=nextline]
\item[\code{Codes.lean}]
Pairing, bounded lists, zig-zag integers, rationals, exponent vectors, sparse
polynomials, rational functions, permutations, finite tables, and generic
proof-tree codes.  Each operation is accompanied by a \code{Primrec} theorem.

\item[\code{SparsePolynomial.lean}]
Canonical sorting and coefficient collection, evaluation semantics, and the
fundamental theorem
\[
 \operatorname{eval}(\operatorname{pnorm}(p),\rho)
 =\operatorname{eval}(p,\rho).
\]
Prove the list-loop invariants first in ordinary Lean; then reflect them into
\code{PRAAx} using the existing primitive-recursion infrastructure.

\item[\code{RationalFunction.lean}]
Nonzero-denominator records, cross-multiplication equality, variable renaming,
and substitution.  This module proves the exact identities for the five
symmetric generators and the generic product polynomial.

\item[\code{LowDegreeCertificates.lean}]
Codes for the quadratic, Cardano, biquadratic, and Ferrari witness packages;
root-list constructors; and Boolean factorization checkers.  The semantic
correctness lemmas can be connected to the existing factorization theorems,
while the object-theory results use only the sparse-normalizer equations.

\item[\code{FinitePermutationGroups.lean}]
Bounded enumeration of \code{Equiv.Perm (Fin d)}, table construction,
commutators, subgroup closure, derived series, and
\code{solvableB}.  Include compact certificates for
$[S_5,S_5]=A_5$ and $[A_5,A_5]=A_5$ rather than relying only on a large
\code{native_decide} computation.

\item[\code{RadicalTowerCertificate.lean}]
The proof-carrying finite-extension and radical-layer records, local checker,
normal-layer compiler, exact-sequence certificates, and the theorem that every
accepted tower produces an accepted solvable-series certificate.

\item[\code{GenericQuinticCertificate.lean}]
Codes for $e_1,\ldots,e_5$ and $P$; permutation renaming; the map from a
permutation to an automorphism table; the inverse map reading the images of
$x_i$; and the checked isomorphism with \code{S5Code}.

\item[\code{AbelRuffiniPRA.lean}]
Definitions of \code{RadSplit}, \code{UniversalRadicalFormula}, and the
generic specialization compiler.  Combine the tower-solvability theorem,
generic Galois certificate, and $S_5$ computation into the ordinary iPRA
derivations.

\item[\code{AbelRuffiniNative.lean}]
Translate the ordinary formulas to the selected native calculus and discharge
them with \code{native_iff_pra}.  This file should contain no algebra.

\item[\code{Audit.lean}]
\texttt{\#check} and \texttt{\#print axioms} for every public endpoint, plus closed
examples for the degree-two, degree-three, degree-four, and $S_5$ checkers.
\end{description}

\subsection{Suggested Lean interfaces}

The exact first-order syntax names should follow ConciseIPRA, but the public
mathematical interfaces should look like the following.  These declarations
are schematic: the archive's \code{Lean/IPRAAbelRuffiniBlueprint.lean} records
the module boundaries without pretending that the missing proofs have already
been checked.

\begin{lstlisting}[style=lean,language={},caption={Certificate-level endpoints}]
structure LowDegreeWitness (n : Fin 5) where
  coeffs : List QCode
  data   : ExprCode
  proof  : ProofCode

def lowCheck (n : Fin 5) : LowDegreeWitness n -> Bool
def lowRoots (n : Fin 5) : LowDegreeWitness n -> List ExprCode

theorem lowCheck_sound {n} {w : LowDegreeWitness n} :
    lowCheck n w = true -> factorCheck w.coeffs (lowRoots n w) = true

def towerCheck : TowerCode -> Bool
def solvableB : GroupCode -> Bool

theorem towerCheck_implies_solvable {t} :
    towerCheck t = true -> solvableB (splittingGroup t) = true

theorem genericQuintic_galCode :
    genericQuinticGroup = S5Code

theorem S5_not_solvableB : solvableB S5Code = false

theorem no_generic_radical_split :
    forall t, RadicalSplit genericQuintic t -> False
\end{lstlisting}

The object-theory endpoint then has the shape

\begin{lstlisting}[style=lean,language={},caption={Object-theory endpoints, schematic}]
theorem pra_lowDegree_factorization (n : Fin 5) :
  [] |-PRA- lowDegreeSentence n

theorem pra_abelRuffini_degree_five :
  [] |-PRA- not (existsCode (UniversalRadicalFormula 5))

theorem pra_abelRuffini_degree_ge_five :
  [] |-PRA- forallNat fun n =>
    5 <= n -> not (existsCode (UniversalRadicalFormula n))

theorem native_abelRuffini_degree_ge_five :
  NPrf [] (NFm.ofPRA abelRuffiniSentence)
\end{lstlisting}

Here \code{|-PRA-} abbreviates ProveIt's actual notation
\(\Gamma \vdash[\mathrm{PRAAx}] A\).  The final proof is a one-line use of
\code{native_iff_pra} after the ordinary derivation exists.

\subsection{How to internalize a primitive-recursive checker}

For each checker, use the same four-stage pattern.

\begin{enumerate}
\item Define a total Lean function on natural-number codes.  Malformed input
  returns \code{false} or a fixed error code.
\item Prove the function primitive recursive compositionally.  Do not mark it
  recursive by an opaque kernel shortcut.
\item State its loop invariant as an equality between primitive-recursive
  functions.  Prove that invariant in \code{PRAAx} by quantifier-free
  induction, using the equations for primitive recursion already available in
  ConciseIPRA.
\item Derive the externally readable semantic theorem from the accepted
  certificate.  The semantic theorem may use structures and fields; the
  object-theory theorem mentions only the code and checker equations.
\end{enumerate}

For fixed identities such as the Cardano and Ferrari factorizations, an even
smaller route is available: generate a closed proof trace of sparse-ring
normalization and verify that trace with a primitive-recursive step checker.
The iPRA proof then establishes soundness of the generic step checker once and
applies it to a concrete numeral coding the trace.

\subsection{Milestones and completion criterion}

A clear path to completion is:

\begin{enumerate}
\item \textbf{Low-degree milestone.}  Kernel-check
  \code{pra_lowDegree_factorization} for $n=1,2,3,4$.  This validates the
  coding and normalization layer before any Galois theory is internalized.
\item \textbf{Finite-group milestone.}  Kernel-check
  \code{S5_not_solvableB} with no native oracle in the public proof term.
\item \textbf{Generic-action milestone.}  Kernel-check the mutually inverse
  maps between the $120$ renamings and the automorphism certificate.
\item \textbf{Radical-tower milestone.}  Kernel-check the compiler from an
  accepted radical tower to a solvable-series certificate.
\item \textbf{Ordinary iPRA endpoint.}  Kernel-check the degree-five
  contradiction and the padding theorem in \code{PRAAx}.
\item \textbf{Native endpoint.}  Transfer through
  \code{native_iff_pra}, run \texttt{\#print axioms}, and place the result under
  the repository's existing no-placeholder audit.
\end{enumerate}

The formalization is complete only when the last two endpoints print no
\code{sorryAx}, no extra arithmetic axiom, and no unrecorded computation
oracle.  The Python/SymPy script in this archive is useful regression evidence
but is not part of that trust claim.
